\section{The deconvolution at level one half}\label{es:sec:level}

\subsection{The forward equation}

Write $s=1+iv$ and, for a zero $\rho=\sigma_\rho+i\gamma_\rho$, $a_\rho=1-\sigma_\rho\in(0,1)$. Proposition~\ref{prop:struct} and the Hadamard product give, with the zeros paired as $\rho,\bar\rho$ so that the sum converges,
\begin{equation}\label{es:eq:lorentz}
u_1(v)=\frac1\pi\Re\frac{\xi'}\xi(1+iv)=\frac1\pi\sum_\rho\frac{a_\rho}{a_\rho^2+(v-\gamma_\rho)^2}=\sum_\rho P_{a_\rho}(v-\gamma_\rho),
\end{equation}
a sum of Cauchy densities $P_a(x)=(a/\pi)/(a^2+x^2)$ of widths $a_\rho$ centred at the ordinates. Under the Riemann hypothesis every width is $\tfrac12$ and
\begin{equation}\label{es:eq:forward}
u_1=P_{1/2}*\mathfrak U_{1/2},\qquad\mathfrak U_{1/2}=\sum_\rho\delta_{\gamma_\rho}:
\end{equation}
the level one density is the Poisson integral of the zero counting measure at distance $\tfrac12$ from the line. In general \eqref{es:eq:lorentz} is the Poisson integral of the positive measure $\sum_\rho\delta_{(\gamma_\rho,a_\rho)}$, each atom smoothed at its own distance $a_\rho$, and Theorem~\ref{tf:thm:deconv} shows that a positive solution $\mu$ of $P_a*\mu=u_1$ exists for a given $a\in(0,\tfrac12]$ if and only if no zero has $a_\rho<a$.

\subsection{The Fourier side}

With $\widehat f(\omega)=\int f(v)e^{-i\omega v}dv$, $\widehat{P_a}(\omega)=e^{-a|\omega|}$, so termwise
\begin{equation}\label{es:eq:u1hat}
\widehat{u_1}(\omega)=\sum_\rho e^{-a_\rho|\omega|}e^{-i\omega\gamma_\rho},
\end{equation}
and under the Riemann hypothesis $\widehat{u_1}=e^{-|\omega|/2}\,\widehat{\mathfrak U_{1/2}}$ with $\widehat{\mathfrak U_{1/2}}(\omega)=\sum_\rho e^{-i\omega\gamma_\rho}$, a tempered distribution. On the other hand, from the Dirichlet series, as tempered distributions,
\begin{equation}\label{es:eq:spectrum}
\widehat{u_1}=1+\widehat G-\sum_{n\ge2}\frac{\Lambda(n)}n\bigl(\delta_{\log n}+\delta_{-\log n}\bigr),
\end{equation}
where the $1$ is the transform of $\delta_0$, the boundary value on $\Re s=1$ of $\pi^{-1}\Re(s-1)^{-1}$, which the prime series does not contain, and $G(v)=\pi^{-1}[1/(1+v^2)-\tfrac12\log\pi+\tfrac12\Re\psi(\tfrac{1+iv}2)]$ is the Gamma part, whose transform is the distribution
\begin{equation}\label{es:eq:Ghat}
\langle\widehat G,g\rangle=(1-\log\pi-\gamma_E)\,g(0)+\sum_{k\ge1}\Bigl[\frac{g(0)}{k+1}-\int_{\mathbb R}g(\omega)\,e^{-(2k+1)|\omega|}\,d\omega\Bigr].
\end{equation}
Away from the origin $\widehat G(\omega)=-e^{-3|\omega|}/(1-e^{-2|\omega|})=-1/(2|\omega|)+1+O(|\omega|)$, the transform of the growth $(2\pi)^{-1}\log|v|$ of $G$: a singularity that is not integrable, and \eqref{es:eq:Ghat} is its finite part regularisation. Multiplying \eqref{es:eq:spectrum} by $e^{|\omega|/2}$ multiplies the atom at $\pm\log n$ by $\sqrt n$.

\begin{theorem}[the explicit formula is the deconvolution]\label{es:thm:explicit}
For every $g\in C_c^\infty(\mathbb R)$, with the sum over $\rho$ absolutely convergent,
\begin{multline}\label{es:eq:weil}
\sum_\rho\int_{\mathbb R}g(\omega)\,e^{(\sigma_\rho-\frac12)|\omega|}\,e^{-i\omega\gamma_\rho}\,d\omega=\int_{\mathbb R}g(\omega)\,e^{|\omega|/2}\,d\omega\\
-\sum_{n\ge2}\frac{\Lambda(n)}{\sqrt n}\bigl(g(\log n)+g(-\log n)\bigr)+(1-\log\pi-\gamma_E)\,g(0)\\
+\sum_{k\ge1}\Bigl[\frac{g(0)}{k+1}-\int_{\mathbb R}g(\omega)\,e^{-(2k+\frac12)|\omega|}\,d\omega\Bigr].
\end{multline}
For even $g$ this is Theorem~\ref{dc:thm:weil}, the explicit formula of Guinand \cite{Guinand48} and Weil \cite{Weil52}; for odd $g$ both sides vanish. Under the Riemann hypothesis the left side is $\sum_\rho\widehat g(\gamma_\rho)$, and \eqref{es:eq:weil} says that on $C_c^\infty(\mathbb R)$
\[
\sum_\rho e^{-i\omega\gamma_\rho}=e^{|\omega|/2}\bigl(1+\widehat G(\omega)\bigr)-\sum_{n\ge2}\frac{\Lambda(n)}{\sqrt n}\bigl(\delta_{\log n}+\delta_{-\log n}\bigr),
\]
where $e^{|\omega|/2}\widehat G$ is the functional $g\mapsto\langle\widehat G,e^{|\cdot|/2}g\rangle$ given by \eqref{es:eq:Ghat}. The left side is tempered; the right side is tempered only as a whole, since the prime measure has mass about $2e^{X/2}$ on $[0,X]$.
\end{theorem}

\begin{proof}
(a) For $\sigma>1$, $\xi'/\xi(s)=1/s+1/(s-1)-\tfrac12\log\pi+\tfrac12\psi(s/2)+\zeta'/\zeta(s)$. Let $\sigma\downarrow1$ with $s=\sigma+iv$. The left side tends to $\pi u_1$ in $\mathcal S'(\mathbb R)$, since $\xi'/\xi$ is analytic across $\Re s=1$ and grows like $\log|v|$. On the right, $\Re1/s\to1/(1+v^2)$, $\tfrac12\Re\psi(s/2)\to\tfrac12\Re\psi(\tfrac{1+iv}2)$, $\pi^{-1}\Re(s-1)^{-1}=P_{\sigma-1}(v)\to\delta_0$, and $\Re\zeta'/\zeta(s)=-\sum_n\Lambda(n)n^{-\sigma}\cos(v\log n)$, whose transform $-\pi\sum_n\Lambda(n)n^{-\sigma}(\delta_{\log n}+\delta_{-\log n})$ converges in $\mathcal S'$ as $\sigma\downarrow1$ by dominated convergence, the measure $\sum_n\Lambda(n)n^{-1}\delta_{\log n}$ having mass $O(X)$ on $[0,X]$. This gives \eqref{es:eq:spectrum}. The digamma series gives $G(v)=(1-\log\pi-\gamma_E)/2\pi+\sum_{k\ge1}[1/(2\pi(k+1))-P_{2k+1}(v)]$, the $1/(1+v^2)$ cancelling the term $k=0$, and transforming termwise gives \eqref{es:eq:Ghat}; the series converges because $\int g(\omega)e^{-b|\omega|}d\omega=2g(0)/b+O(b^{-2})$ and $1/(k+1)-2/(2k+1)=O(k^{-2})$.

(b) By \eqref{es:eq:lorentz} the partial sums of $\sum_\rho P_{a_\rho}(v-\gamma_\rho)$ are positive and increase to $u_1=O(\log|v|)$, so the series converges in $\mathcal S'$ and $\widehat{u_1}=\sum_\rho e^{-a_\rho|\omega|}e^{-i\omega\gamma_\rho}$.

(c) Let $g\in C_c^\infty(\mathbb R)$ and $\varphi=e^{|\omega|/2}g$, which is Lipschitz with compact support and has $\varphi'$ of bounded variation. Let $\varphi_\epsilon=\varphi*\eta_\epsilon$ with a mollifier $\eta_\epsilon$: then $\varphi_\epsilon\in C_c^\infty$ with supports in a fixed compact set, $\varphi_\epsilon\to\varphi$ uniformly, and $\|\varphi_\epsilon'\|_\infty$ and the total variation of $\varphi_\epsilon'$ are bounded uniformly in $\epsilon$. Pair (a) and (b) with $\varphi_\epsilon$. On the zero side the term of $\rho$ is the transform of $e^{-a_\rho|\omega|}\varphi_\epsilon$ at $\gamma_\rho$; two integrations by parts, the second against the measure $d\bigl(e^{-a_\rho|\omega|}\varphi_\epsilon\bigr)'$ of bounded mass, bound it by $C(1+\gamma_\rho^2)^{-1}$ uniformly in $\epsilon$ and $\rho$, and $\sum_\rho(1+\gamma_\rho^2)^{-1}<\infty$, so by dominated convergence the zero side tends to $\sum_\rho\int g(\omega)e^{(1/2-a_\rho)|\omega|}e^{-i\omega\gamma_\rho}d\omega$, the left side of \eqref{es:eq:weil}, since $\tfrac12-a_\rho=\sigma_\rho-\tfrac12$. On the other side the prime sum is finite on the support, $\int\varphi_\epsilon\to\int g\,e^{|\omega|/2}$, and $\langle\widehat G,\varphi_\epsilon\rangle\to\langle\widehat G,\varphi\rangle$, because each term of \eqref{es:eq:Ghat} converges and $|\varphi_\epsilon(\omega)-\varphi_\epsilon(0)|\le C|\omega|$ bounds the $k$-th term by $C/k^2$ uniformly. This is \eqref{es:eq:weil}. Under the Riemann hypothesis $a_\rho=\tfrac12$, and $N(T)=O(T\log T)$ makes $\sum_\rho e^{-i\omega\gamma_\rho}$ tempered.
\end{proof}

The theorem is classical, and that is the point, made already in Section~\ref{dc:sec:weil}: the backward Cauchy flow from level $1$ to level $\tfrac12$, the object that Chapters~\ref{chap:lo} and~\ref{chap:cx} computed by EM, is the explicit formula, and the statement of Theorem~\ref{tf:thm:deconv} that the Thorin measure at level $\theta$ is the solution of $P_{1-\theta}*\mu=u_1$ is, at $\theta=\tfrac12$ and under the Riemann hypothesis, the statement that the zeros are the Fourier transform of the primes. The multiplier $e^{|\omega|/2}$ of the deconvolution, whose unboundedness made the problem ill posed, is the passage from $\Lambda(n)/n$ to $\Lambda(n)/\sqrt n$, from a tempered measure, of mass about $X$ on $[0,X]$, to one of mass about $2e^{X/2}$, which pairs only with test functions of compact support. Everything that was said about the conditioning of the grid matrix in Chapter~\ref{chap:cx} is a statement about this passage.

\subsection{What the ill posedness costs}

\begin{proposition}[an unresolvable perturbation]\label{es:prop:resolution}
Let $\Omega>0$, $L\ge1$, $\chi_L(v)=(2\pi L^2)^{-1/2}e^{-v^2/2L^2}$, $N=\int_{\mathbb R}|\cos(\Omega v)|\,\chi_L(v)\,dv$ and $\nu=N^{-1}\cos(\Omega v)\,\chi_L(v)\,dv$. Then $\|\nu\|_{TV}=1$, $\widehat\nu$ is concentrated within $O(1/L)$ of $\pm\Omega$, and
\[
\|P_{1/2}*\nu\|_\infty\le\sqrt{\pi/8}\;L^{-1}e^{1/8L^2}\,e^{-\Omega/2}\le0.72\,e^{-\Omega/2}.
\]
Consequently, if $0.72\,e^{-\Omega/2}\le2\varepsilon$ and $\mu$ is a positive measure with $\mu\ge N^{-1}\chi_L\,dv$, then $\mu$ and $\mu+\nu$ are positive measures at total variation distance $1$ whose data $P_{1/2}*\mu$ and $P_{1/2}*(\mu+\nu)$ lie within $\varepsilon$ in $L^\infty$ of one function: data of accuracy $\varepsilon$ do not distinguish them.
\end{proposition}

\begin{proof}
$\widehat{\chi_L}(\omega)=e^{-L^2\omega^2/2}$, so $\widehat\nu(\omega)=(2N)^{-1}\bigl(e^{-L^2(\omega-\Omega)^2/2}+e^{-L^2(\omega+\Omega)^2/2}\bigr)$. By Fourier inversion, the evenness of $e^{-|\omega|/2}$ and $e^{-|\omega|/2}\le e^{-\omega/2}$,
\[
\begin{split}\|P_{1/2}*\nu\|_\infty&\le\frac1{2\pi}\int_{\mathbb R}e^{-|\omega|/2}|\widehat\nu(\omega)|\,d\omega\\&\le\frac1{2\pi N}\int_{\mathbb R}e^{-\omega/2}e^{-L^2(\omega-\Omega)^2/2}\,d\omega=\frac{e^{1/8L^2}}{\sqrt{2\pi}\,L\,N}\,e^{-\Omega/2}.\end{split}
\]
From $|\cos x|=\frac2\pi+\frac4\pi\sum_{k\ge1}(-1)^{k+1}\frac{\cos2kx}{4k^2-1}$ and $\int\cos(2k\Omega v)\chi_L(v)\,dv=e^{-2k^2L^2\Omega^2}$,
\[
N=\frac2\pi+\frac4\pi\sum_{k\ge1}(-1)^{k+1}\frac{e^{-2k^2L^2\Omega^2}}{4k^2-1}\ge\frac2\pi,
\]
the series being alternating with decreasing terms. With $N\ge2/\pi$ this is the bound, and $\sqrt{\pi/8}\,e^{1/8}=0.7101$. For the last statement, $\nu\ge-N^{-1}\chi_L$, so $\mu+\nu\ge0$, and the midpoint of the two data lies within $\tfrac12\|P_{1/2}*\nu\|_\infty\le\varepsilon$ of each.
\end{proof}

\begin{remark}[heuristic]\label{es:rem:resolution}
Identifying resolution at scale $\delta$ with frequency $\Omega=\pi/\delta$, Proposition~\ref{es:prop:resolution} suggests that a general density cannot be resolved at scale $\delta$ from data less accurate than about $\tfrac12e^{-\pi/2\delta}$, and, since in \eqref{es:eq:weil} a test function supported in $[-\Omega,\Omega]$ involves $\Lambda(n)$ only for $n\le e^\Omega$, that the primes to $e^{\pi/\delta}$ are the matching arithmetic input. Neither is a theorem: no loss function is fixed, no method is excluded, and both fail for measures with few atoms, as the computations below show.
\end{remark}

For $\delta=0.1$ the heuristic gives $\varepsilon\approx10^{-7}$ and primes to $e^{31}\approx10^{13}$; for $\delta=0.01$ it is $3\times10^{-69}$ and $10^{136}$. The EM computation of Chapter~\ref{chap:cx} reached $\delta\approx10^{-3}$ in the positions of well separated atoms from data accurate to $10^{-20}$, which is better than the heuristic allows for a general measure, and the reason is that the target is not a general measure. Proposition~\ref{es:prop:resolution} is about the worst perturbation; a positive measure that is a sum of a few atoms per unit length has far fewer degrees of freedom than a density, and the question is whether a method can use that.

\section{Sparse recovery}\label{es:sec:sparse}

\subsection{Positivity and sparsity}

Two facts from the theory of super resolution apply directly. First, Cand\`es and Fernandez Granda \cite{CandesFG14} proved that a measure $\mu=\sum_kc_k\delta_{t_k}$ on $[0,1]$ is recovered exactly from its Fourier coefficients $\widehat\mu(\omega)$, $|\omega|\le\Omega$, by minimising the total variation subject to the data, provided the atoms are separated by at least $2\cdot2\pi/\Omega$; with noise $\varepsilon$ the recovery error is $O(\varepsilon)$ times a super resolution factor, polynomial in $\Omega\delta$ rather than exponential \cite{CandesFG13}. Second, for positive measures the separation is not needed at all: de Castro and Gamboa \cite{deCastroGamboa12}, and Schiebinger, Robeva and Recht \cite{SchiebingerRR18}, showed that a positive measure with $k$ atoms is the unique positive measure with its moments of orders $0$ to $2k$, and is recovered by a linear program. Our measure is positive, and on a window of length $W$ at height $v$ it has about $W\log(v/2\pi)/2\pi$ atoms separated on average by $2\pi/\log(v/2\pi)$, with occasional pairs much closer. So the problem $P_{1/2}*\mu=u_1$, $\mu\ge0$, restricted to a window and to the frequencies the data support, is of the kind these theorems treat.

The gridded convex methods of Chapter~\ref{chap:cx} already respect positivity, and the sparse vertex solutions of nonnegative least squares are a step toward sparsity; but on a grid of spacing $h$ the model has $W/h$ degrees of freedom, the unknown is a density, and the conditioning is that of Proposition~\ref{es:prop:resolution}. The gain from sparsity comes only when the model is the set of sums of atoms, and for that the natural tool is on the Fourier side.

\subsection{Lorentzians become exponentials}

\begin{theorem}[tapered transform]\label{es:thm:taper}
Let $w(v)=e^{-(v-c)^2/2\ell^2}$, $\ell>0$, and
\[
F(\omega)=\int_{\mathbb R}u_1(v)\,w(v)\,e^{i\omega(v-c)}\,dv .
\]
Then for $\omega\ell^2>1$
\begin{equation}\label{es:eq:taperF}
F(\omega)=\sum_\rho w(\gamma_\rho+ia_\rho)\,e^{-a_\rho\omega}e^{i\omega(\gamma_\rho-c)}+R(\omega),
\end{equation}
where $w(x+iy)=w(x)e^{y^2/2\ell^2}e^{-iy(x-c)/\ell^2}$ is the entire extension of the taper, the series converges absolutely, and, with $Y=\omega\ell^2$,
\[
\begin{gathered}|R(\omega)|\le\frac1\pi\,e^{-\ell^2\omega^2/2}\sum_\rho a_\rho\int_{\mathbb R}\frac{e^{-(x-c)^2/2\ell^2}\,dx}{(x-\gamma_\rho)^2+(Y-a_\rho)^2},\\ \int_{\mathbb R}\frac{e^{-(x-c)^2/2\ell^2}\,dx}{(x-\gamma)^2+B^2}\le\min\Bigl(\frac\pi B,\frac{\sqrt{2\pi}\,\ell}{B^2}\Bigr),\end{gathered}
\]
so that, heuristically, $R$ is of the order $e^{-\ell^2\omega^2/2}$ times the number of zeros per unit length at height $c$, divided by $\omega\ell$. In particular, for $\omega_j=\omega_0+j\Delta$ with $\omega_0\ell^2>1$,
\begin{gather*}
F(\omega_j)=\sum_\rho A_\rho\,z_\rho^{\,j}+R(\omega_j),\qquad|R(\omega_j)|\le C_0\,e^{-\ell^2\omega_j^2/2},\\
A_\rho=w(\gamma_\rho+ia_\rho)e^{-a_\rho\omega_0}e^{i\omega_0(\gamma_\rho-c)},\qquad z_\rho=e^{\Delta(-a_\rho+i(\gamma_\rho-c))},
\end{gather*}
where $C_0$ is the bound on $|R(\omega_0)|$ divided by $e^{-\ell^2\omega_0^2/2}$. The Gamma factor does not enter: \eqref{es:eq:lorentz} is all of $u_1$, and the Lorentzians $P_{2k+1}$ centred at $0$, which the trivial zeros $-2k$ contribute to $\Re\zeta'/\zeta$, are cancelled in $\xi'/\xi$ by the poles of $\Gamma(s/2)$.
\end{theorem}

\begin{proof}
By \eqref{es:eq:lorentz}, $u_1=\sum_\rho P_{a_\rho}(\cdot-\gamma_\rho)$ exactly, with positive terms, and $\int u_1w<\infty$; so $F(\omega)=\sum_\rho\int P_{a_\rho}(v-\gamma_\rho)\phi(v)\,dv$ with $\phi(v)=w(v)e^{i\omega(v-c)}$, which is entire with
\[
|\phi(x+iy)|=e^{-(x-c)^2/2\ell^2}\,e^{y^2/2\ell^2-\omega y}.
\]
Write $P_a(v-\gamma)=(2\pi i)^{-1}[(v-\gamma-ia)^{-1}-(v-\gamma+ia)^{-1}]$ and move the integral over $\mathbb R$ to the line $\Im v=Y$, $Y=\omega\ell^2$, where $|\phi|$ is smallest. Since $0<a<1<Y$, the first term picks up the residue at $v=\gamma+ia$, the second has no pole in between, the vertical sides contribute nothing by the Gaussian decay in $x$, and
\[
\int_{\mathbb R}P_a(v-\gamma)\phi(v)\,dv=\phi(\gamma+ia)+\frac1{2\pi i}\int_{\Im v=Y}\frac{2ia\,\phi(v)}{(v-\gamma)^2+a^2}\,dv .
\]
On the line, $|\phi|=e^{-(x-c)^2/2\ell^2}e^{-\ell^2\omega^2/2}$ and $|(v-\gamma)^2+a^2|=|v-\gamma-ia|\,|v-\gamma+ia|\ge(x-\gamma)^2+(Y-a)^2$, which gives the term of the bound attached to $\rho$; the two bounds for the integral follow from $\int((x-\gamma)^2+B^2)^{-1}dx=\pi/B$ and $\int e^{-(x-c)^2/2\ell^2}dx=\sqrt{2\pi}\,\ell$. The sums over $\rho$ converge: $|w(\gamma+ia)|\le e^{1/2\ell^2}e^{-(\gamma-c)^2/2\ell^2}$, the integrals are $O((\gamma-c)^{-2})$ uniformly because $Y-a\ge Y-1>0$, and $N(T)=O(T\log T)$. Finally $\phi(\gamma+ia)=w(\gamma+ia)e^{i\omega(\gamma-c)}e^{-a\omega}$, and the bound at $\omega_j$ holds because the bound on $|R(\omega)|$ divided by $e^{-\ell^2\omega^2/2}$ decreases in $\omega$.
\end{proof}

A hypothesis of the form $\omega\ell^2>a_\rho$ is needed, and $\omega\ell^2>1$ suffices: for $\omega\ell^2\le a_\rho$ the line $\Im v=\omega\ell^2$ does not reach the pole, and the residue is not picked up. The settings of the computations below ($\omega_0\ell^2=500$ at height $200$, $20$ at height $3\times10^{12}$) satisfy it.

The content of the theorem is in the last display. The samples of the tapered transform are a finite sum of damped complex exponentials in $j$, one per zero that the taper sees, plus a remainder, with the pole $z_\rho$ carrying both coordinates of the zero: its modulus gives $a_\rho=1-\sigma_\rho$, its argument gives $\gamma_\rho-c$ modulo $2\pi/\Delta$. Nothing is assumed about the real parts in the statement; equation \eqref{es:eq:taperF} is \eqref{es:eq:lorentz} and nothing else.

\subsection{Prony and the matrix pencil}

Recovering the poles of a sum of exponentials from $J$ equally spaced samples is Prony's problem, and the matrix pencil of Hua and Sarkar \cite{HuaSarkar90} solves it stably. Form the Hankel matrix $H_{ik}=F(\omega_{i+k})$, $0\le i<J-L$, $0\le k<L$, with $L\approx J/2$, and its two submatrices $H_0$, $H_1$ obtained by deleting the last and the first row. If the samples are an exact sum of $M$ exponentials with distinct nonzero poles and nonzero amplitudes, and $M\le\min(L,J-L-1)$, then $H_0$ and $H_1$ have rank $M$ and the values of $z$ at which $H_1-zH_0$ drops rank are the poles $z_\rho$ (Lemma~\ref{es:lem:pencil}). With noise, one takes the rank $M$ truncated singular value decomposition $H_0\approx U_M\Sigma_MV_M^*$, chooses $M$ by the gap in the singular values, and computes the eigenvalues of $U_M^*H_1V_M\Sigma_M^{-1}$. The amplitudes $A_\rho$ follow by least squares against the Vandermonde matrix of the recovered poles, and a component is accepted as a zero when the modulus of its amplitude matches the prediction $|w(\gamma_\rho+ia_\rho)|e^{-a_\rho\omega_0}$ of Theorem~\ref{es:thm:taper}, a diagnostic and not a certificate (in the runs at levels $\theta<1$ below, the prediction is evaluated at the width $\theta-\tfrac12$ of a zero on the line); spurious components, which absorb the zeros at the edge of the taper and the wrapped poles of zeros farther away, have amplitudes orders of magnitude off the profile.

\begin{lemma}[exact matrix pencil]\label{es:lem:pencil}
Let $f_j=\sum_{m=1}^MA_mz_m^{\,j}$, $0\le j<J$, with distinct nonzero $z_m$, nonzero $A_m$ and $M\le\min(L,J-L-1)$, and let $H_0=(f_{i+k})$ and $H_1=(f_{i+k+1})$, $0\le i\le J-L-2$, $0\le k<L$. Then $\operatorname{rank}H_0=\operatorname{rank}H_1=M$, $\operatorname{rank}(H_1-zH_0)<M$ exactly when $z\in\{z_1,\dots,z_M\}$, and if $H_0=U_M\Sigma_MV_M^*$ is a reduced singular value decomposition then $U_M^*H_1V_M\Sigma_M^{-1}$ is similar to $\operatorname{diag}(z_1,\dots,z_M)$.
\end{lemma}

\begin{proof}
With $V=(z_m^{\,i})$ of size $(J-L-1)\times M$, $W=(z_m^{\,k})$ of size $L\times M$, $D=\operatorname{diag}(A_m)$ and $Z=\operatorname{diag}(z_m)$, we have $H_0=VDW^{T}$ and $H_1=VDZW^{T}$. $V$ and $W$ are Vandermonde with distinct nodes and at least $M$ rows, so both have rank $M$; hence $\operatorname{rank}H_0=\operatorname{rank}H_1=M$ and $\operatorname{rank}(H_1-zH_0)=\operatorname{rank}(Z-zI)=M-\#\{m:z_m=z\}$. The columns of $U_M$ span the column space of $V$, so $P=U_M^*V$ is invertible; from $H_0V_M=U_M\Sigma_M$ we get $PD(W^TV_M)=\Sigma_M$, so $W^TV_M=D^{-1}P^{-1}\Sigma_M$ and $U_M^*H_1V_M\Sigma_M^{-1}=PDZD^{-1}P^{-1}=PZP^{-1}$.
\end{proof}

Heuristically, the error analysis is the standard one for the pencil: if the samples carry an error $\eta$ and the exponentials are separated, the poles move by $O(\eta/|A_{\min}|)$ times a condition number polynomial in $J$ and in the separation, not exponential in the bandwidth. The bandwidth enters through $|w(\gamma_\rho+ia_\rho)|e^{-a_\rho\omega_{J-1}}$: at the largest frequency used, the signal of a zero is $e^{-\Omega/2}$ times its amplitude, and the data error at that frequency is $\eta$, so the effective noise is $\eta e^{\Omega/2}$. This is the same $e^{\Omega/2}$ as in Proposition~\ref{es:prop:resolution}, but it now multiplies the data error, not the resolution: with data exact to $10^{-20}$ and $\Omega=45$ the effective noise is $10^{-10}$, and with double precision transforms, $\eta=10^{-16}$, it is $10^{-6}$ at the top frequency and far smaller at the frequencies that carry most of the signal.

\section{The computation}\label{es:sec:computation}

The numbers of this section are the author's, computed by the scripts of the source and not rerun for this book. The data are $u_1$ on $[150,250]$ at spacing $0.02$, enclosed in balls of radius below $5\times10^{-21}$ by the script of Chapter~\ref{chap:cx}; the taper is centred at $c=200$ with $\ell=10$, so that the zeros within $15$ of the centre have taper weight at least $0.3$ and those beyond $30$ weight below $0.012$. The transform $F(\omega_j)$ is computed as a Riemann sum over the data in ball arithmetic, at $\omega_j=5+0.1j$, $j<400$, so $\Omega=45$, and converted to double precision. The balls enclose the finite sum, not $F$: the aliasing error of the sum is of the order $e^{-a(2\pi/h-\Omega)}=e^{-135}$ with $a=\tfrac12$ and $h=0.02$, but the truncation of the data at $c\pm5\ell$, where the taper is $e^{-12.5}=3.7\times10^{-6}$, is not bounded. The Hankel matrix is $200\times200$; the singular values fall from $1$ to $10^{-9}$ over the first sixty and the model order is taken there. At $\Delta=0.1$ the ordinates are determined modulo $2\pi/\Delta=62.8$, so zeros farther than $31.4$ from the centre appear wrapped into the window, with taper weights below $e^{-4.9}$; the amplitude test sets them aside, and only ordinates within $25$ of the centre are reported. The whole computation takes two minutes, almost all in the ball arithmetic transform; in double precision it takes a second.

\begin{table}[ht]
\centering
\small
\begin{tabular}{lcccccc}
\toprule
& zeros & \multicolumn{2}{c}{position error} & \multicolumn{2}{c}{error in $\sigma_\rho$}\\
$|\gamma-c|\le$ & & max & median & max & median\\
\midrule
$10$ & 11 & $4.5\times10^{-10}$ & $1.8\times10^{-10}$ & $7.9\times10^{-10}$ & $2.3\times10^{-10}$\\
$15$ & 17 & $5.0\times10^{-7}$ & $2.5\times10^{-10}$ & $8.4\times10^{-7}$ & $3.8\times10^{-10}$\\
$20$ & 21 & $5.0\times10^{-7}$ & $4.5\times10^{-10}$ & $8.4\times10^{-7}$ & $7.9\times10^{-10}$\\
$25$ & 28 & $4.1\times10^{-5}$ & $2.7\times10^{-9}$ & $4.0\times10^{-5}$ & $1.1\times10^{-8}$\\
\bottomrule
\end{tabular}
\caption{The matrix pencil at height $200$, level one data, taper $\ell=10$: errors in the position and in the real part of the recovered zeros, by distance from the centre (computed by the author's scripts, not rerun).}\label{es:tab:height200}
\end{table}

\noindent The eleven zeros within $10$ of the centre are placed to within $8\times10^{-10}$ in both coordinates; the error grows toward the edge of the taper as the amplitudes fall, as the heuristic error analysis predicts. For comparison, the backward flow of Chapter~\ref{chap:cx} on the same data, with $4000$ accelerated EM steps and the certificate, placed the centroids to about $10^{-3}$. The gain is seven orders of magnitude at the centre, from the same $u_1$. The modulus of the tapered transform $F(\omega)$ follows the line $e^{-\omega/2}$, so that the signal falls by $e^{-20}$ over the frequencies used (the author's figure is not reproduced).

The real parts are the new element. Every recovered pole has $|z_\rho|=e^{-\Delta a_\rho}$ with $a_\rho=0.5000000$ to the digits shown; the fit does not constrain the real parts, and would have reported otherwise. To make the point we added to $u_1$ the Lorentzians of a fictitious pair of zeros at $0.85\pm60.089i$ and $0.15\pm60.089i$, which is what a zero off the line and its reflection would contribute, and ran the pencil on a window centred at $60$. It returned two components at $\gamma=60.08900000$ and $60.08899905$ with $a=0.15000000$ and $0.84999991$, beside the genuine zeros $59.347044$ and $60.831779$ with $a=0.500000$. The pencil reads the real parts of the zeros off the level one data, to seven digits, without any assumption.

\section{Posterior concentration, and what it cannot do}\label{es:sec:posterior}

One may ask for a posterior concentration result for the recovery of $\mathfrak U_{1/2}$, and hence for a proof of the Riemann hypothesis. The first can be done under hypotheses; the second cannot, and the theorem itself shows why.

Fix a window, centre $c$ and taper $\ell$, a sampling step $\Delta$, the samples $\omega_j=\omega_0+j\Delta$, $0\le j<J$, with $\omega_0\ell^2>1$, and a radius $R$ with $2R\Delta<2\pi$, so that $\gamma\mapsto e^{i\Delta(\gamma-c)}$ is one to one on $[c-R,c+R]$. By Theorem~\ref{es:thm:taper} the data vector $F=(F(\omega_j))_j$ is
\[
F=\Phi(\vartheta^0)+b,\qquad\Phi(\vartheta)_j=\sum_{m=1}^MA_m\,z_m^{\,j},\qquad z_m=e^{\Delta(-a_m+i(\gamma_m-c))},
\]
where $\vartheta^0$ collects the zeros with $|\gamma_m-c|\le R$ and $b$ is the deterministic bias from the zeros beyond $R$ and the remainder of Theorem~\ref{es:thm:taper}; its size is about $2n(c)(\ell^2/R)e^{-R^2/2\ell^2}$ per sample, with $n(c)$ the density of zeros at height $c$. To make a Bayesian statement we model the observation as $Y=F+\tau\xi$, where $\xi\in\mathbb C^J$ has independent standard normal real and imaginary parts, and put a prior on $\vartheta$.

\begin{theorem}[posterior concentration on a window]\label{es:thm:concentration}
Fix $s>0$, $0<a_-<a_+<1$, $0<\alpha_-<A_+$ and an integer $M_{\max}$ with $J\ge2M_{\max}$. For $0\le M\le M_{\max}$ let $K_M$ be the compact set of $\vartheta=(\gamma_m,a_m,A_m)_{m\le M}$ with $c-R\le\gamma_1<\dots<\gamma_M\le c+R$, $\gamma_{m+1}-\gamma_m\ge s$, $a_m\in[a_-,a_+]$ and $\alpha_-\le|A_m|\le A_+$. Let the prior be $\Pi=\sum_Mp_M\Pi_M$, where $\Pi_M$ has a density on $K_M$ with values between $\pi_->0$ and $\pi_+$. Let $\vartheta^0$ be an interior point of $K_{M^0}$ with $p_{M^0}>0$, and let $Y=\Phi(\vartheta^0)+b+\tau\xi$ with $\|b\|_2\le B\tau$. Put $d(\vartheta,\vartheta^0)=\max_m\bigl(|\gamma_m-\gamma^0_m|+|a_m-a^0_m|+|A_m-A^0_m|\bigr)$ if $M=M^0$ and $d=\infty$ otherwise. Then there are $\tau_0,r_0,K,\kappa_0>0$, depending on $\vartheta^0$, the sets $K_M$, the prior, $J$, $\Delta$ and $B$, such that for $\tau\le\tau_0$ and $K\tau\le\varepsilon\le r_0$, with probability at least $1-e^{-\kappa_0\varepsilon^2/\tau^2}$,
\[
\Pi\bigl(d(\vartheta,\vartheta^0)\le\varepsilon\bigm|Y\bigr)\ \ge\ 1-e^{-\kappa_0\varepsilon^2/\tau^2}.
\]
In particular, for fixed $\varepsilon_1>0$ and data with $\|b\|_2\le B\tau$, as $\tau\to0$ the posterior probability that every component has $|\sigma_m-\tfrac12|\le\varepsilon_1$ tends to $1$ in probability if every zero of $\vartheta^0$ has $|\sigma^0_m-\tfrac12|<\varepsilon_1$, and to $0$ if one has $|\sigma^0_m-\tfrac12|>\varepsilon_1$.
\end{theorem}

\begin{proof}
\emph{Identifiability.} On each $K_M$ the poles $z_m$ are distinct and nonzero, and $(\gamma,a)\mapsto z$ is one to one on $[c-R,c+R]\times(0,1)$ because $2R\Delta<2\pi$. If $\Phi(\vartheta)=\Phi(\vartheta^0)$ with $\vartheta\in K_M$, the difference is a combination of at most $M+M^0\le J$ distinct exponentials vanishing at $J$ consecutive indices; the Vandermonde matrix of the distinct poles has full column rank, so after merging equal poles every coefficient vanishes, and since all amplitudes are nonzero each pole of $\vartheta$ is a pole of $\vartheta^0$ with the same amplitude and conversely. Hence $M=M^0$ and, the labels being fixed by the order of the $\gamma_m$, $\vartheta=\vartheta^0$. By continuity and compactness $g=\min_{M\ne M^0}\min_{K_M}\|\Phi(\vartheta)-\Phi(\vartheta^0)\|_2>0$. In the $4M^0$ real coordinates the derivative of $\Phi$ at $\vartheta^0$ is injective: with respect to $(z_m,A_m)$ it is the $J\times2M^0$ matrix with columns $(z^0_m)^j$ and $jA^0_m(z^0_m)^{j-1}$, whose first $2M^0$ rows have determinant $\pm\prod_mA^0_m\prod_{m<n}(z^0_n-z^0_m)^4\ne0$, and $dz=z\Delta(-da+i\,d\gamma)$ with $z\ne0$. So $\|\Phi(\vartheta)-\Phi(\vartheta^0)\|_2\ge\lambda_0d$ for $d\le r_1$, and on the compact set $\{d\ge r_1\}\cap K_{M^0}$ the left side has a positive minimum; hence $\|\Phi(\vartheta)-\Phi(\vartheta^0)\|_2\ge\lambda d(\vartheta,\vartheta^0)$ on $K_{M^0}$ for some $\lambda>0$, and $\|\Phi(\vartheta)-\Phi(\vartheta^0)\|_2\le L_\Phi d$ on the $d$-ball of radius $r_1$ about $\vartheta^0$, which lies in $K_{M^0}$ if $r_1$ is small.

\emph{Noise.} Put $u=\Phi(\vartheta)-\Phi(\vartheta^0)$ and $e=b+\tau\xi$. The log likelihood ratio of $\vartheta$ to $\vartheta^0$ is $\bigl(2\Re\langle u,e\rangle-\|u\|^2\bigr)/2\tau^2$. On the event $E_t=\{\|\xi\|_2\le\sqrt{2J}+\sqrt{2t}\}$, of probability at least $1-e^{-t}$ by Gaussian concentration, $\|e\|\le q\tau$ with $q=B+\sqrt{2J}+\sqrt{2t}$, so the ratio lies between $e^{-\|u\|^2/2\tau^2-q\|u\|/\tau}$ and $e^{-\|u\|^2/4\tau^2+q^2}$, by $-x^2/2+qx\le-x^2/4+q^2$.

\emph{Denominator.} On $\{M=M^0,\ d\le\tau\}$, $\tau\le r_1$, we have $\|u\|\le L_\Phi\tau$, and this set has prior mass at least $p_{M^0}\pi_-v_0\tau^D$, $D=4M^0$, where $v_0\tau^D$ is the volume of the $d$-ball of radius $\tau$. So the integral of the ratio against $\Pi$ is at least $p_{M^0}\pi_-v_0\tau^De^{-L_\Phi^2/2-L_\Phi q}$.

\emph{Numerator.} Since $e^{-\lambda^2d^2/4\tau^2}\le e^{-\lambda^2\varepsilon^2/8\tau^2}e^{-\lambda^2d^2/8\tau^2}$ for $d>\varepsilon$, the integral of the ratio over $\{M=M^0,\ d>\varepsilon\}$ is at most
\[
p_{M^0}\pi_+e^{q^2}e^{-\lambda^2\varepsilon^2/8\tau^2}\int_{\mathbb R^D}e^{-\lambda^2d(\vartheta,\vartheta^0)^2/8\tau^2}\,d\vartheta=p_{M^0}\pi_+C_D\Bigl(\frac\tau\lambda\Bigr)^De^{q^2}e^{-\lambda^2\varepsilon^2/8\tau^2},
\]
and the integral over $\{M\ne M^0\}$ is at most $e^{q^2}e^{-g^2/4\tau^2}$.

\emph{Conclusion.} The powers $\tau^D$ cancel:
\[
\Pi(d>\varepsilon\mid Y)\le C_1e^{q^2+L_\Phi q}\bigl(e^{-\lambda^2\varepsilon^2/8\tau^2}+\tau^{-D}e^{-g^2/4\tau^2}\bigr).
\]
For $\varepsilon\le r_0:=g/\lambda$ and $\tau\le\tau_0$, with $\tau_0\le r_1$ so small that $\tau^{-D}e^{-g^2/8\tau^2}\le1$, the second term is at most the first. Since $q^2+L_\Phi q\le2q^2+L_\Phi^2\le C_2+12t$, on $E_t$ we get $\Pi(d>\varepsilon\mid Y)\le C_3e^{12t}e^{-\lambda^2\varepsilon^2/8\tau^2}$. Take $t=\lambda^2\varepsilon^2/192\tau^2$: then $\Pi(d>\varepsilon\mid Y)\le C_3e^{-\lambda^2\varepsilon^2/16\tau^2}\le e^{-\lambda^2\varepsilon^2/32\tau^2}$ once $\varepsilon\ge K\tau$ with $K^2\lambda^2\ge32\log\max(C_3,1)$, on an event of probability at least $1-e^{-\lambda^2\varepsilon^2/192\tau^2}$. This is the claim with $\kappa_0=\lambda^2/192$. For the last statement, $|\sigma_m-\sigma^0_m|=|a_m-a^0_m|\le d$; apply the bound with the fixed radius $\varepsilon=\varepsilon_1-\max_m|\sigma^0_m-\tfrac12|$ in the first case and $\varepsilon=\tfrac12(|\sigma^0_m-\tfrac12|-\varepsilon_1)$ in the second, reduced to $r_0$ if necessary.
\end{proof}

\begin{remark}\label{es:rem:settings}
The theorem needs the bias below the noise, $\|b\|_2\le B\tau$ with $B$ of order one, and no aliasing, $R<\pi/\Delta$. At the settings of the computations of this chapter these are incompatible. With $\ell=10$, $\Delta=0.1$ and $\tau=10^{-16}$ at height $200$, the bias falls below $\tau$ only for $R\ge86$, while $\pi/\Delta=31.4$; with $\ell=2$ and $\Delta=0.2$ at height $3\times10^{12}$ it needs $R\ge17.4$, while $\pi/\Delta=15.7$. With $R$ just below $\pi/\Delta$ the bias is about $2\times10^{-3}$ per sample at height $200$, and the theorem then holds only with $B$ of the order of $\|b\|_2/\tau$, that is at radii of the order of the bias divided by $\lambda$, which is no information. The theorem describes a computation with $\Delta<\pi/86=0.036$ at $\ell=10$, about $1100$ samples on $[5,45]$, or one in which the amplitudes are tied to $(\gamma_m,a_m)$ through the taper, which removes the aliasing; that variant is not proved here.

The separation $s>0$ excludes two components at one ordinate, and so the pairs $\tfrac12\pm\delta$ that the functional equation imposes on zeros off the line: such a configuration is not an interior point of any $K_M$, and the rate $\tau$ of the theorem does not apply to it. Near the line the two poles of such a pair merge; Section~\ref{es:sec:singular} treats that singular case.
\end{remark}

Under its hypotheses the theorem is a useful statement: the posterior for the zeros in a window, including their real parts, contracts to the truth at the rate of the data accuracy. It is the Bayesian form of Theorem~\ref{es:thm:taper}. Its hypotheses do not hold at the settings of the computations above (Remark~\ref{es:rem:settings}): the errors of $10^{-10}$ to $10^{-7}$ that the matrix pencil reaches from data with $\tau=10^{-16}$ are an observation, not an instance of the theorem. Its constants, however, depend on the separation of the poles, and Section~\ref{es:sec:singular} shows that this dependence is not a technicality: the functional equation forces off line zeros to come in pairs whose poles merge as they approach the line, and near the line the rate for the real part of a pair is $\tau^{1/2}$, not $\tau$ (Theorem~\ref{es:thm:singular}).

The theorem concentrates the posterior at $\vartheta^0$, whatever $\vartheta^0$ is. It says that the data decide, in the limit of exact data, the real parts of finitely many zeros to any finite precision; it does not say what they decide. If the Riemann hypothesis is true, the posterior probability that the zeros in the window lie within $\varepsilon$ of the line tends to one; if it is false for some zero in the window, the same probability tends to zero; and the theorem is indifferent between the two cases, because posterior consistency is a statement about the mechanism, not about the parameter. To prove the hypothesis from it one would need, first, $\varepsilon=0$, which a posterior at a positive noise level does not give, and second, all windows at once, which no finite computation covers. The second is the obstacle that the verification of the hypothesis to height $3\times10^{12}$ meets: a finite computation settles finitely many zeros, and the hypothesis is a statement about all of them. The first is particular to the posterior: sign changes of $Z$ together with a count of zeros prove $\sigma=\tfrac12$ exactly for every zero they cover, while a posterior at a positive noise level gives the real parts only to finite precision. Nor is there a loophole in the limit $\tau\to0$: the data $u_1$ are not random, the noise is a modelling device, and the limit of the posterior is the point mass at the true zeros, which is where the problem started. We record this because posterior concentration is exactly the right description of what the pencil computation does, but the step from concentration to a theorem about the parameter is the step that no statistical argument can make: consistency means the posterior ends at the truth, and the truth here is the Riemann hypothesis or its negation, undetermined by the theorem.

What the theorem does add to the programme is a precise sense in which $\mathfrak U_{1/2}$ is recoverable: under its hypotheses the atomic model is identifiable from the level one data on a window without aliasing, the posterior contracts at the rate of the data accuracy, and the real parts are part of what is recovered. That is the sparse answer to the deconvolution, and it is as far as the data go.
